% GENERATED FILE. Edit canonical lessons, then run npm run generate:books.
% canonical-source-hash: fnv1a64:abd5353b41e89582
% canonical-lessons: KA-C317-bandare, KA-C317-male-bandare, KA-R317-if-and-otherwise

\chapter{If It Rains}
\label{ch:ka-if-it-rains}
\hlchaptermodality{317}

\begin{chapteropening}
\textbf{By the end of this chapter:} \emph{I can say what will happen if something happens with the -are form (bandare, iddare), add ondu vēḷe for if by chance, and say if not with -adiddare.}

Complete the last lesson of chapter 317: say what will happen if, if by chance and if not, with -are and -adiddare.
\end{chapteropening}

\section[\texorpdfstring{\kn{ಬಂದರೆ} (bandare)}{ಬಂದರೆ (bandare)}]{\kn{ಬಂದರೆ} (bandare) — if}

\label{lesson:KA-C317-bandare}



\begin{quote}
Say \textquotedblleft{}but\textquotedblright{} (\emph{ādare}), \textquotedblleft{}in that case\textquotedblright{} (\emph{hāgādare}) and \textquotedblleft{}otherwise\textquotedblright{} (\emph{illadiddare}). Then \textquotedblleft{}I have been learning Kannada for two months.\textquotedblright{} (\emph{nānu eraḍu tiṅgaḷinda kannaḍa kaliyuttiddēne}.)

Those three words share an ending. This lesson takes it out.
\end{quote}

\subsection*{You'll want to know: the -\kn{ಅರೆ} form}
\begin{quote}
\textbf{\kn{ನೀನು} \kn{ಬಂದರೆ} \kn{ನಾವು} \kn{ಒಟ್ಟಿಗೆ} \kn{ಊಟ} \kn{ಮಾಡೋಣ}.}
\end{quote}

\begin{quote}
\emph{nīnu bandare nāvu oṭṭige ūṭa māḍōṇa.} — \textquotedblleft{}If you come, let's eat together.\textquotedblright{}
\end{quote}

\textbf{\kn{ಬಂದರೆ}} is the past without its ending, \textbf{\kn{ಬಂದ}-}, with \textbf{-\kn{ಅರೆ}}: \textquotedblleft{}if (you) come\textquotedblright{}. It stays the same for every person: \textbf{\kn{ನಾನು} \kn{ಬಂದರೆ}}, \textbf{\kn{ಅವಳು} \kn{ಬಂದರೆ}}, \textbf{\kn{ನಾವು} \kn{ಬಂದರೆ}}. Though it is built on the past, it points ahead: \textquotedblleft{}if you come (tomorrow)\textquotedblright{}. The same build on other verbs: \textbf{\kn{ಮಾಡಿದ}-} $\to$ \textbf{\kn{ಮಾಡಿದರೆ}}, \textbf{\kn{ಹೋದ}-} $\to$ \textbf{\kn{ಹೋದರೆ}}, \textbf{\kn{ಇದ್ದ}-} $\to$ \textbf{\kn{ಇದ್ದರೆ}}, \textbf{\kn{ಆದ}-} $\to$ \textbf{\kn{ಆದರೆ}}.

\textbf{\kn{ಮಾಡೋಣ}}, \textquotedblleft{}let's do\textquotedblright{}, has the ending of \textbf{\kn{ನಾಳೆ} \kn{ಸಿಗೋಣ}}.

Now the three words. \textbf{\kn{ಆದರೆ}} is \textquotedblleft{}if it becomes so\textquotedblright{}, and Kannada uses it for \textquotedblleft{}but\textquotedblright{}. \textbf{\kn{ಹಾಗಾದರೆ}} is \textbf{\kn{ಹಾಗೆ}} + \textbf{\kn{ಆದರೆ}}, \textquotedblleft{}if it is that way\textquotedblright{}: \textquotedblleft{}in that case\textquotedblright{}. \textbf{\kn{ಇಲ್ಲದಿದ್ದರೆ}} is \textquotedblleft{}if it is not\textquotedblright{}: \textquotedblleft{}otherwise\textquotedblright{}. Even \textbf{\kn{ಯಾಕೆಂದರೆ}} holds one: \textquotedblleft{}if (you) say why\textquotedblright{}, so \textquotedblleft{}because\textquotedblright{}.

\subsection*{You'll want to know: \kn{ಒಂದು} \kn{ವೇಳೆ}, if by chance}
\begin{quote}
\textbf{\kn{ಒಂದು} \kn{ವೇಳೆ} \kn{ಬಸ್} \kn{ತಡವಾದರೆ}, \kn{ನನಗೆ} \kn{ಫೋನ್} \kn{ಮಾಡಿ}.}
\end{quote}

\begin{quote}
\emph{ondu vēḷe bas taḍavādare, nanage phōn māḍi.} — \textquotedblleft{}If by chance the bus is late, phone me.\textquotedblright{}
\end{quote}

\textbf{\kn{ಒಂದು} \kn{ವೇಳೆ}}, \textquotedblleft{}in case, if by chance\textquotedblright{}, opens the clause; \textbf{-\kn{ಅರೆ}} still closes it. \textbf{\kn{ತಡವಾದರೆ}} is \textbf{\kn{ತಡ}} (\emph{taḍa}, \textquotedblleft{}delay\textquotedblright{}) + \textbf{\kn{ಆದರೆ}}: \textquotedblleft{}if it becomes late\textquotedblright{}. \textbf{\kn{ಒಂದು} \kn{ವೇಳೆ}} says you do not expect it, but you are ready for it.

\subsection*{Your turn}
Say these aloud:

\begin{itemize}
\raggedright
  \item \textquotedblleft{}If you come, let's eat together\textquotedblright{} — \emph{nīnu bandare nāvu oṭṭige ūṭa māḍōṇa}
  \item \textquotedblleft{}If she comes, I will tell her\textquotedblright{} — \emph{avaḷu bandare nānu hēḷuttēne}
  \item \textquotedblleft{}If by chance the bus is late, phone me\textquotedblright{} — \emph{ondu vēḷe bas taḍavādare, nanage phōn māḍi}
\end{itemize}

\begin{tcolorbox}[breakable,colback=teal!4,colframe=teal!35!black,title={Before you move on}]
\textquotedblleft{}If (you) come\textquotedblright{}? (\textbf{\kn{ಬಂದರೆ}}.) \textquotedblleft{}If (you) go\textquotedblright{}? (\textbf{\kn{ಹೋದರೆ}}.) \textquotedblleft{}If there is\textquotedblright{}? (\textbf{\kn{ಇದ್ದರೆ}}.) Which of the words you knew is \textquotedblleft{}if it becomes\textquotedblright{}? (\textbf{\kn{ಆದರೆ}}, \textquotedblleft{}but\textquotedblright{}.) \textquotedblleft{}If by chance\textquotedblright{}? (\textbf{\kn{ಒಂದು} \kn{ವೇಳೆ}} … \textbf{-\kn{ಅರೆ}}.)
\end{tcolorbox}
\section[\texorpdfstring{\kn{ಮಳೆ} \kn{ಬಂದರೆ} (maḷe bandare)}{ಮಳೆ ಬಂದರೆ (maḷe bandare)}]{\kn{ಮಳೆ} \kn{ಬಂದರೆ} (maḷe bandare) — if it rains}

\label{lesson:KA-C317-male-bandare}



\begin{quote}
Say \textquotedblleft{}If you come, let's eat together.\textquotedblright{} (\emph{nīnu bandare nāvu oṭṭige ūṭa māḍōṇa}.) Then \textquotedblleft{}otherwise\textquotedblright{}. (\emph{illadiddare}.)
\end{quote}

\subsection*{You'll want to know: if this, then that}
\begin{quote}
\textbf{\kn{ನಾಳೆ} \kn{ಮಳೆ} \kn{ಬಂದರೆ} \kn{ನಾನು} \kn{ಮನೆಯಲ್ಲೇ} \kn{ಇರುತ್ತೇನೆ}.}
\end{quote}

\begin{quote}
\emph{nāḷe maḷe bandare nānu maneyallē iruttēne.} — \textquotedblleft{}If it rains tomorrow, I will stay at home.\textquotedblright{}
\end{quote}

The if-clause comes first and the result after it, in the \textbf{-\kn{ಉತ್ತೇನೆ}} that says \textquotedblleft{}will\textquotedblright{} as well as \textquotedblleft{}do\textquotedblright{}. For something that holds every time, the same shape:

\begin{quote}
\textbf{\kn{ಬಿಸಿಲು} \kn{ಇದ್ದರೆ} \kn{ನಾವು} \kn{ಕಡಲತೀರಕ್ಕೆ} \kn{ಹೋಗುತ್ತೇವೆ}.}
\end{quote}

\begin{quote}
\emph{bisilu iddare nāvu kaḍalatīrakke hōguttēve.} — \textquotedblleft{}When it is sunny, we go to the beach.\textquotedblright{}
\end{quote}

\textbf{\kn{ಹೋಗುತ್ತೇವೆ}} (\emph{hōguttēve}, \textquotedblleft{}we go\textquotedblright{}), the \textquotedblleft{}we\textquotedblright{} of \textbf{\kn{ಹೋಗುತ್ತೇನೆ}}.

\subsection*{You'll want to know: if not}
\begin{quote}
\textbf{\kn{ನೀನು} \kn{ಬೇಗ} \kn{ಬರದಿದ್ದರೆ} \kn{ಬಸ್} \kn{ಹೋಗುತ್ತದೆ}.}
\end{quote}

\begin{quote}
\emph{nīnu bēga baradiddare bas hōguttade.} — \textquotedblleft{}If you don't come quickly, the bus will go.\textquotedblright{}
\end{quote}

\textbf{\kn{ಬರದಿದ್ದರೆ}} is \textbf{\kn{ಬರದೆ}} (\emph{barade}, \textquotedblleft{}without coming\textquotedblright{}) with \textbf{\kn{ಇದ್ದರೆ}}: \textquotedblleft{}if (you) are without coming\textquotedblright{}. It is built the way \textbf{\kn{ಇಲ್ಲದಿದ್ದರೆ}} is, \textbf{\kn{ಇಲ್ಲದೆ}} + \textbf{\kn{ಇದ್ದರೆ}}. More of them: \textbf{\kn{ಮಾಡದಿದ್ದರೆ}}, \textquotedblleft{}if (you) don't do it\textquotedblright{}; \textbf{\kn{ಹೋಗದಿದ್ದರೆ}}, \textquotedblleft{}if (you) don't go\textquotedblright{}.

\subsection*{Your turn}
Say these aloud:

\begin{itemize}
\raggedright
  \item \textquotedblleft{}If it rains tomorrow, I will stay at home\textquotedblright{} — \emph{nāḷe maḷe bandare nānu maneyallē iruttēne}
  \item \textquotedblleft{}When it is sunny, we go to the beach\textquotedblright{} — \emph{bisilu iddare nāvu kaḍalatīrakke hōguttēve}
  \item \textquotedblleft{}If you don't come quickly, the bus will go\textquotedblright{} — \emph{nīnu bēga baradiddare bas hōguttade}
  \item one thing you will do if it rains tomorrow, and one if it does not
\end{itemize}

\begin{tcolorbox}[breakable,colback=teal!4,colframe=teal!35!black,title={Before you move on}]
\textquotedblleft{}If it rains\textquotedblright{}? (\textbf{\kn{ಮಳೆ} \kn{ಬಂದರೆ}}.) Which form says what will happen? (\textbf{-\kn{ಉತ್ತೇನೆ}}, the form that says \textquotedblleft{}do\textquotedblright{} and \textquotedblleft{}will\textquotedblright{}.) \textquotedblleft{}If you don't come\textquotedblright{}? (\textbf{\kn{ಬರದಿದ್ದರೆ}}.) Which word you knew is built the same way? (\textbf{\kn{ಇಲ್ಲದಿದ್ದರೆ}}.)
\end{tcolorbox}
\section[(review)]{Review — if, and otherwise}

\label{lesson:KA-R317-if-and-otherwise}



\begin{quote}
Say \textquotedblleft{}a medicine\textquotedblright{} and \textquotedblleft{}clean\textquotedblright{}. (\emph{auṣadha}, \emph{svaccha}.) Which letter opens the first, and which sits small inside the second? (\textbf{\kn{ಔ}}; \textbf{\kn{ಛ}}.)

\begin{itemize}
\raggedright
  \item \emph{Write it:} \textbf{\kn{ಔಷಧ}} and \textbf{\kn{ಸ್ವಚ್ಛ}}, once each from memory
\end{itemize}
\end{quote}

\subsection*{Your turn}
Say these aloud:

\begin{itemize}
\raggedright
  \item If you come tomorrow, let's go to the cinema.
  \item If by chance it rains, phone me.
  \item If it rains, I will stay at home. Otherwise I will come at six.
  \item If you don't take this medicine, you must go to the hospital.
  \item If the train is faster than the bus, we must go by train.
  \item This is the smallest room. If the guests come, I will sleep here.
  \item If you want to become a doctor, you must study.
\end{itemize}

\begin{tcolorbox}[breakable,colback=teal!4,colframe=teal!35!black,title={Before you move on}]
\textquotedblleft{}If (you) come\textquotedblright{}? (\textbf{\kn{ಬಂದರೆ}}.) \textquotedblleft{}If by chance\textquotedblright{}? (\textbf{\kn{ಒಂದು} \kn{ವೇಳೆ}} … \textbf{-\kn{ಅರೆ}}.) \textquotedblleft{}If you don't come\textquotedblright{}? (\textbf{\kn{ಬರದಿದ್ದರೆ}}.) \textquotedblleft{}Otherwise\textquotedblright{}? (\textbf{\kn{ಇಲ್ಲದಿದ್ದರೆ}}.)
\end{tcolorbox}
